\subsection*{通信与隐私}
真人交流测试准备发送者、收件人和第三方三个账号，逐步执行投递、接受和发送消息。虚拟互动测试则检查待回复、回复成功、失败重试及屏蔽状态。两类流程各自使用独立夹具，便于定位权限错误和状态错误。

\begin{table}[H]\centering\small
\caption{通信与私有状态用例}
\begin{tabularx}{\linewidth}{@{}p{15mm}p{38mm}X@{}}\toprule
编号 & 前置与操作 & 预期结果及检查点 \\ \midrule
IT-09 & 收件人接受已收藏对象纸条；第三方也尝试 & 本人建立唯一会话和首轮消息；第三方不能处理 \\
IT-10 & 双方发送消息；第三方访问同一会话 & 双方记录和时间可读；第三方越权被拒绝 \\
IT-11 & 虚拟回复失败后重试；伪造内部补写动作 & 原用户消息不重复；外部不能调用内部完成路径 \\
IT-12 & 记录心笺、变更授权、推进周期与重置 & 分数不隐含意愿，周期重新判断；私人备注不进入上游 \\ \bottomrule
\end{tabularx}\end{table}

\DesignCode{real-chat.test.mjs}验证纸条到会话的完整链路与第三方隔离；\DesignCode{core.test.mjs}、\DesignCode{actions.test.mjs}、\DesignCode{state.test.mjs}分别检查规则、动作白名单和持久化。错误输入除返回失败外，还应保持原状态和记录数不变。

\DesignCode{server.test.mjs}截获发送给模拟模型的消息，检查其中不含心笺备注和客户端伪造的历史。测试在请求中主动加入这些字段，再核对上游实际收到的内容，从而检查服务端是否按规定组织模型素材。

重置测试确认虚拟纸条被清空，账号仍可使用。交流参考的查询测试确认双方开启时结果可见，任一方关闭后再次查询时结果隐藏。
